% Trabalho acadêmico nas normas da ABNT com a classe abnTeX2 (https://www.abntex.net.br/).
% Troque os dados abaixo e escreva os capítulos em capitulos/.
\documentclass[
  12pt,
  openany,       % capítulos começam em qualquer página (use openright para impressão frente e verso)
  oneside,       % troque por twoside para imprimir frente e verso
  a4paper,
  english,       % idioma adicional (abstract)
  brazil         % idioma principal: o último da lista
]{abntex2}

\usepackage{fontspec}
\usepackage{indentfirst}         % recua o primeiro parágrafo de cada seção
\usepackage{graphicx}
\usepackage{microtype}
\usepackage[brazilian,hyperpageref]{backref}
\usepackage[alf]{abntex2cite}    % citações autor-data (alf) ou numéricas (num)

% páginas citadas na bibliografia
\renewcommand{\backrefpagesname}{Citado na(s) página(s):~}
\renewcommand{\backref}{}
\renewcommand*{\backrefalt}[4]{%
  \ifcase #1 Nenhuma citação no texto.%
  \or Citado na página #2.%
  \else Citado #1 vezes nas páginas #2.%
  \fi}

% ---------------------------------------------------------------- dados do trabalho
\titulo{Título do trabalho}
\autor{Nome do Autor}
\local{Cidade}
\data{2026}
\orientador{Prof. Dr. Nome do Orientador}
\instituicao{%
  Nome da Instituição
  \par
  Nome do Departamento ou Curso}
\tipotrabalho{Trabalho de Conclusão de Curso}
\preambulo{Trabalho de Conclusão de Curso apresentado ao curso de Nome do Curso da Nome da Instituição, como requisito parcial para a obtenção do título de Bacharel.}

% links e informações do PDF
\makeatletter
\hypersetup{
  pdftitle={\@title},
  pdfauthor={\@author},
  pdfsubject={\imprimirpreambulo},
  colorlinks=true,
  linkcolor=black,
  citecolor=black,
  urlcolor=blue,
  bookmarksdepth=4
}
\makeatother

\setlength{\parindent}{1.3cm}
\setlength{\parskip}{0.2cm}

\makeindex

\begin{document}

\selectlanguage{brazil}
\frenchspacing

% ---------------------------------------------------------------- elementos pré-textuais
\imprimircapa
\imprimirfolhaderosto

\begin{resumo}
  Escreva aqui o resumo do trabalho, em um único parágrafo de 150 a 500 palavras, com o
  objetivo, o método, os resultados e as conclusões.

  \textbf{Palavras-chave}: palavra-chave. outra palavra-chave. mais uma.
\end{resumo}

\begin{resumo}[Abstract]
  \begin{otherlanguage*}{english}
    Write the abstract here.

    \vspace{\onelineskip}
    \noindent
    \textbf{Keywords}: keyword. another keyword.
  \end{otherlanguage*}
\end{resumo}

\pdfbookmark[0]{\listfigurename}{lof}
\listoffigures*
\cleardoublepage

\pdfbookmark[0]{\listtablename}{lot}
\listoftables*
\cleardoublepage

\pdfbookmark[0]{\contentsname}{toc}
\tableofcontents*
\cleardoublepage

% ---------------------------------------------------------------- elementos textuais
\textual

\chapter{Introdução}
\label{cap:introducao}

Este modelo segue as normas da ABNT para trabalhos acadêmicos (NBR 14724), referências
(NBR 6023) e citações (NBR 10520), usando a classe abnTeX2.

As citações seguem o sistema autor-data. Uma citação indireta fica entre parênteses no
fim da frase, como esta sobre a estrutura do trabalho, que tem elementos pré-textuais,
textuais e pós-textuais \cite{abnt-nbr14724}. Quando o autor faz parte da frase, use
\verb|\citeonline|: segundo \citeonline{abntex2modelo}, o modelo pode ser adaptado às
regras de cada instituição.

O \autoref{cap:desenvolvimento} mostra exemplos de figura e tabela.

\chapter{Desenvolvimento}
\label{cap:desenvolvimento}

\section{Figuras}

A \autoref{fig:exemplo} mostra como inserir uma figura com título e fonte, como pede a norma.
Coloque as imagens na pasta \texttt{figuras/} ou cole uma imagem direto no editor.

\begin{figure}[htb]
  \centering
  \caption{Título da figura}
  \label{fig:exemplo}
  \fbox{\parbox[c][4cm][c]{0.6\textwidth}{\centering Troque por \texttt{\textbackslash includegraphics}}}
  \legend{Fonte: elaborado pelo autor.}
\end{figure}

\section{Tabelas}

A \autoref{tab:exemplo} segue o mesmo padrão.

\begin{table}[htb]
  \centering
  \ABNTEXfontereduzida
  \caption{Título da tabela}
  \label{tab:exemplo}
  \begin{tabular}{lcc}
    \hline
    Item & Valor & Unidade \\
    \hline
    Primeiro & 10 & cm \\
    Segundo  & 20 & cm \\
    \hline
  \end{tabular}
  \legend{Fonte: elaborado pelo autor.}
\end{table}

\chapter{Conclusão}
\label{cap:conclusao}

Retome o objetivo do trabalho, resuma os resultados e indique trabalhos futuros.


% ---------------------------------------------------------------- elementos pós-textuais
\postextual

\bibliography{referencias}

\end{document}
